\documentclass[a4paper,12pt]{article}
\usepackage[utf8]{inputenc} % Pacote para acentuação
\usepackage[brazil]{babel} % Para quem vai escrever em português brasileiro
\usepackage[lmargin=3cm,tmargin=3cm,rmargin=2cm,bmargin=2cm]{geometry} % Formato que lembra a ABNT
\usepackage{amsmath, amsthm, amsfonts, amssymb, dsfont, mathtools} % Pacotes matemáticos
\usepackage{graphicx}

\begin{document}

\section{Tabela de Notas} % Use section* para um título sem numeração

\begin{center} % Centraliza a tabela no texto [11:10]
\begin{tabular}{|l|c|c|c|} % Define 4 colunas: 1 alinhada à esquerda (l) e 3 centralizadas (c) 
    \hline % Adiciona uma linha horizontal no topo
    \textbf{Matérias} & \textbf{Prova 1} & \textbf{Prova 2} & \textbf{Nota Final} \\ % Títulos das colunas em negrito 
    \hline % Linha horizontal após os títulos 
    Cálculo 1 & 7,5 & 7 & NFC1 \\ % Exemplo de linha 
    \hline % Linha horizontal entre as matérias 
    Álgebra Linear & 8,5 & 6 & NFAL \\
    \hline
    Química 1 & 6,5 & 6 & NFQ1 \\
    \hline
    Introdução às Engenharias & 10 & 9,5 & NFI \\
    \hline
\end{tabular}
\end{center}

\end{document}
